\documentclass{situationist-preprint}

\title{Must the Past Hypothesis Be Explained?}
\msid{TSM-05-b}
\msfamily{TSM-05}
\msdiscipline{Time, Space and Modality}
\mstype{Consequence}
\mskeywords{past hypothesis; explanation; best system; typicality; fine-tuning}

\begin{document}
\maketitle

\begin{abstract}
TSM-05-a argues that time has a direction because the universe began in a low-entropy state, and that this past hypothesis requires no further explanation. We draw the consequence for the demand that it must be explained. No rational requirement is violated by leaving it unexplained; cosmological attempts to explain it remain legitimate but are not owed; and its alleged improbability supports neither design nor a multiverse.
\end{abstract}

\section{The consequence}

By Theorem~4.5 of TSM-05-a, the past hypothesis (PH) requires no further explanation: it is an axiom of the best system and so located by the best theory (Lemmas~4.1 and~4.2), and it is not surprising relative to that theory, since the statistical measure figures there only conditional on PH (Lemma~4.3) and no rival makes PH much more probable (Lemma~4.4). Three things follow.

\begin{proposition}\label{p:three}
(i) A theory that leaves PH unexplained is not thereby defective. (ii) A cosmological explanation of PH would be welcome but is not owed. (iii) The improbability of PH under an unconditional measure is no evidence for any hypothesis that would make PH probable.
\end{proposition}

\begin{proof}
(i) A theory is defective for leaving a fact unexplained only if the fact requires explanation, which by Theorem~4.5 PH does not. (ii) The inquiry is legitimate, but its failure would not count against the best system; any such theory would cite a further special condition of the same kind (SCI-18-b). (iii) The unconditional measure is not part of the best system (Lemma~4.3), and an improbability computed from a measure that has no application is not evidence \citep{Callender2004}.
\end{proof}

Clause (iii) agrees with REL-10-b: the measure that would make the initial state improbable has a role only conditional on PH and cannot assess it. Clause (ii) fits SCI-18-b, on which science cannot explain PH; none is owed. The answer to the title is no. \citet{Price1996} holds that a low-entropy final state would strike us as calling for explanation, so that to exempt the initial one is a double standard. TSM-05-a accepts the symmetry and denies the demand at both boundaries (Objection~1 there): a statement of either boundary macrostate would earn its place in the best system in the same way.

\begin{thebibliography}{9}
\bibitem[Callender(2004)]{Callender2004} Callender, C. (2004). Measures, explanations and the past: Should \enquote{special} initial conditions be explained? \emph{The British Journal for the Philosophy of Science}, 55(2), 195--217.
\bibitem[Price(1996)]{Price1996} Price, H. (1996). \emph{Time's Arrow and Archimedes' Point}. Oxford University Press.
\end{thebibliography}

\end{document}
